\begin{abstract}
As machine-generated text proliferates online, language models increasingly risk training on their own output, creating a feedback loop driving iterative model collapse that erodes diversity and calibration across generations.
While collapse dynamics are increasingly understood, \emph{detecting} onset remains open: iterative training yields short, noisy time series ($T{<}15$) resisting early-warning diagnostics, collapse is a first-order transition precluding critical-slowing-down signals, and no prior work systematically compares detection methods on this problem.
We map the \emph{detectability landscape} (how detection varies across contamination levels, observation windows, and statistical methods) on five model families (345M--6.9B).
A steep onset emerges: permutation Granger testing finds zero significant canary$\to$downstream pairs through 45\% contamination, then jumps to reliable detection (6/8 pairs; 95\% CI $[0.25, 0.92]$) within two percentage points, defining blind, onset, and detection regimes.
Across five methods, null-control false-positive rates span 7\%--94\%, revealing that method choice alone determines whether collapse appears detectable.
Token entropy precedes diversity decline by ${\geq}1$ generation, enabling proof-of-concept intervention reducing peak perplexity by 64\%\footnote{Code and data are available at \url{https://anonymous.4open.science/r/distrib_canary_collapse-3463/}.}.
\end{abstract}
